\scanpage{176}
% archive sheet number: 175
\diagram{\begin{tikzcd}[ampersand replacement=\&,column sep=small,row sep=small]
\text{multiplicatif}\&\overbrace{\hspace{1.8cm}}^{\deleted{\text{\unsure{le}}}\text{séparable}}\&\mathbb Z/\unsure{l|\ell}\mathbb Z\quad\text{fini \unsure{séparable}}\\
\mathbb G_m\arrow[r]\&G\arrow[ur]\arrow[dr]\&\left.\vphantom{\begin{matrix}a\\a\end{matrix}}\right\}\text{fini}\\
\text{\unsure{séparable} connexe}\&\underbrace{\hspace{1.8cm}}_{\text{connexe}}\&\unsure{\mu_p|H_p}\quad\text{fini radiciel}
\end{tikzcd}}
\diagram{\begin{tikzcd}[ampersand replacement=\&,column sep=small,row sep=small]
\text{discret (ordinaire)}\&\overbrace{\hspace{1.8cm}}^{\text{sans }p\text{-torsion}}\&\mathbb Z/l\mathbb Z\arrow[dl]\quad\text{fini multiplicatif}\\
\mathbb Z\&G\arrow[l]\&\left.\vphantom{\begin{matrix}a\\a\end{matrix}}\right\}\text{fini}\\
\deleted{\text{\unsure{séparable}}}\text{ sans torsion}\&\underbrace{\hspace{1.8cm}}_{\deleted{\text{\unsure{seul}}}\text{sans }l\text{-torsion}}\&\mathbb Z/p\mathbb Z\arrow[ul]\quad\text{fini unipotent}
\end{tikzcd}}

\noindent\rule{\linewidth}{.3pt}

\diagram{\begin{tikzcd}[ampersand replacement=\&,column sep=small,row sep=small]
\text{unipotent}\&\overbrace{\hspace{1.8cm}}^{\text{séparable}}\&\mathbb Z/p\mathbb Z\quad\text{fini séparable}\\
\mathbb G_a\arrow[r]\&G\arrow[ur]\arrow[dr]\&\left.\vphantom{\begin{matrix}a\\a\end{matrix}}\right\}\text{fini}\\
\text{\unsure{séparable} connexe}\&\underbrace{\hspace{1.8cm}}_{\text{connexe}}\&\alpha_p\quad\text{fini radiciel}
\end{tikzcd}}
\diagram{\begin{tikzcd}[ampersand replacement=\&,column sep=small,row sep=small]
\text{connexe (= infinitésimal)}\&\overbrace{\hspace{1.8cm}}^{\text{sans torsion unipotente}}\&\mu_p\arrow[dl]\quad\text{fini multiplicatif}\\
\unsure{\mathbb V|V}\&G\arrow[l]\&\left.\vphantom{\begin{matrix}a\\a\end{matrix}}\right\}\text{fini}\\
\deleted{\text{\unsure{séparable}}}\text{ sans torsion}\&\underbrace{\hspace{1.8cm}}_{\text{sans torsion multiplicative}}\&\alpha_p\arrow[ul]\quad\text{fini unipotent}
\end{tikzcd}}

\uline{N.B.} Si on prend des catégories des \deleted{\unsure{ind.}} groupes \inserted{formels} \unsure{des|ou} presque algébriques, \unsure{sur} un corps alg. clos, les multiplicatifs et l'additif se mélangent (dans le fini)

\margin{torsion $=$ \deleted{plus petit ss-groupe \unsure{tel que le} quotient \unsure{ait} \illegible{2} séparable $=$ plus petit ss groupe} ss-groupe fini}

\[
\left\{\begin{aligned}
\text{multiplicatif}&\; -\;\text{discret ordinaire}\\
\text{unipotent}&\; -\;\text{purement infinitésimal}
\end{aligned}\right.
\]
\[
\left\{\begin{aligned}
\text{connexe séparable}&\; -\;\deleted{\text{séparable}}\ \text{sans torsion}\\
\text{fini séparable}&\; -\;\text{fini multiplicatif}\\
\text{fini radiciel}&\; -\;\text{fini unipotent}
\end{aligned}\right.
\]
\[
\left\{\begin{aligned}
\begin{gathered}\deleted{\text{fini séparable}}\\\text{connexe}\end{gathered}&\; -\;\text{sans torsion multiplicative}\\
\text{séparable}&\; -\;\text{sans torsion additive}\\
\text{fini}&\; =\;\text{fini}
\end{aligned}\right.
\]
\[
\left\{\begin{aligned}
\text{multiplicatif connexe séparable}&\; -\;\text{groupe discret sans torsion}\\
\text{unipotent connexe séparable}&\; -\;\text{groupe infinitésimal \deleted{séparable} sans torsion}\\
\begin{gathered}\text{fini séparable mult.}\\\deleted{\text{connexe séparable}}\end{gathered}&\; =\;\text{fini sép. multipl.}\\
\text{fini radiciel mult}&\; =\;\text{fini sép. unipotent}\\
\text{fini radiciel unipotent}&\; =\;\text{fini radiciel unipotent}
\end{aligned}\right.
\]
% slanted
\margin{\deleted{\unsure{formel} ($=$ \unsure{proalgébrique})}\quad\deleted{(\unsure{au fini radiciel})}}
